\section{Overview}\label{sec:overview}
Margining collateral that has no continuous market price is a problem traditional finance has already solved for bilateral derivatives and trading-book capital, in two components. The International Swaps and Derivatives Association's Standard Initial Margin Model (ISDA SIMM) prescribes sensitivity-based initial margin from published risk weights and correlations \cite{isda_simm_v2_6}. The Basel Fundamental Review of the Trading Book (FRTB) market-risk standard prescribes liquidity horizons and conservative floors where market data does not exist \cite{bcbs_frtb_d457}. D-SIMM combines the two into limits for lending markets, and this paper develops it for tokenised funds, whose shares exit through issuer redemption at the published valuation, the per-share valuation the issuer or its administrator publishes on a stated frequency.

This paper specifies the model. The public companion \emph{Applying D-SIMM} contains the data, parameter-assignment, validation, and implementation material.

A lending market has two sides. Suppliers deposit a tokenised asset as collateral, and borrowers draw cash against it. The over-collateralisation the market demands, one minus the loan-to-value ratio, plays the role initial margin plays on a derivative: it reserves against the move in collateral value before the position can be closed. When a position crosses its liquidation threshold, the deployment, the lending market instance holding the collateral, recovers its cash through the issuer's redemption process. Each of the four outputs governs one part of the risk of that recovery, for each collateral, $a$.

The borrow limit $\mathrm{LTV}^{(a)}$ is the loan-to-value ratio at which no further borrowing is permitted. The liquidation threshold $\mathrm{LLTV}^{(a)}$ (liquidation loan-to-value) is the ratio at which a position becomes eligible for liquidation. The liquidation buffer between the two gives a borrower room before recovery starts. The borrow cap $B_{\max}^{(a)}$ bounds the debt outstanding against the collateral, so that one dealing window clears it. The supply cap $S_{\max}^{(a)}$ bounds the collateral the market accepts, so that the whole portfolio clears within a small number of dealing windows. In what follows we write $\theta^{(a)}:=\mathrm{LLTV}^{(a)}$ and $\ell_{\max}^{(a)}:=\mathrm{LTV}^{(a)}$, so that ratios such as $\ell_{\max}^{(a)}/\theta^{(a)}$ express proximity to liquidation in normalised units. Figure~\ref{fig:architecture} outlines the framework: attested terms and published risk weights feed the D-SIMM haircut engine, and the haircut sets the on-chain limits. Due diligence admits or rejects an asset, and monitoring can only reduce the limits for new exposure.

% ==================
% Figure: architecture of D-SIMM.
%
% Composition: the calculation runs left to right across the top, unbroken.
% The two things that govern it from outside sit underneath its two ends and
% point up into it, the due-diligence screen at the entry and the monitor at
% the on-chain output. Nothing else is drawn.
%
% Spacing: 4mm grid. 2mm box padding (half unit), 12mm vertical separation
% between the calculation and what governs it (3 units), 16mm between stages
% (4 units), calculation boxes 12mm tall (3 units), the other two 8mm (2 units).
%
% Ink: one accent for the calculation, one neutral grey for what sits outside
% it. Solid for the single pass, dashed for a later valuation cycle. One stroke
% weight, one arrowhead. Natural size, no scaling, so the figure's type matches
% the body type. For a monochrome figure, set dsAccent to black and dsMuted to
% black!58; every distinction above survives it.
% ==================
\begin{figure}[ht]
\centering
\definecolor{dsAccent}{HTML}{15546F}
\definecolor{dsMuted}{HTML}{4A5157}
\begin{tikzpicture}[
  font=\small,
  >={Stealth[length=1.5mm,width=1.2mm]},
  line width=0.55pt,
  every node/.style={inner xsep=2.2mm, inner ysep=1.6mm, align=left},
  pass/.style={draw=dsAccent, fill=dsAccent!8, minimum height=12mm},
  aside/.style={draw=dsMuted, minimum height=8mm},
  flow/.style={->, draw=dsAccent},
  back/.style={->, draw=dsMuted, dash pattern=on 1.3mm off 1.0mm},
  lab/.style={font=\footnotesize, text=dsMuted, inner xsep=1.1mm, inner ysep=0.8mm},
]
  % the pass, left to right
  \node[pass]                    (in)  {Attested terms\\Published risk weights};
  \node[pass, right=16mm of in]  (eng) {D-SIMM haircut engine};
  \node[pass, right=16mm of eng] (lim) {On-chain limits\\$\theta,\ \ell_{\max}$, caps};

  % what governs the pass from outside it
  \node[aside, below=12mm of in] (dd)  {Due diligence};
  \node[aside] (mon) at (dd -| lim)    {Monitoring};

  \draw[flow] (in)  -- (eng);
  \draw[flow] (eng) -- node[lab, above] {$H$} (lim);
  \draw[flow] (dd)  -- node[lab, right] {admit or reject} (in);
  \draw[back] (mon) -- node[lab, left]  {reduces limits} (lim);
\end{tikzpicture}
\caption{Attested terms and published risk weights feed the D-SIMM haircut
engine. The haircut $H$ sets the active on-chain limits with no iteration. Due
diligence admits or rejects an asset, and monitoring can only reduce those limits.}
\label{fig:architecture}
\end{figure}


\subsection{Universe and settlement}\label{sec:overview-channel}
A tokenised asset is a token whose ownership is governed by existing legal frameworks. Some have physical backing held off-chain, such as gold and share certificates. Assets without physical backing may in time be governed entirely by on-chain ownership rights, derivatives and equities among them. Few such assets exist today. The collateral universe is tokenised assets that correspond to an underlying fund. The risk factors every type is charged against are fixed in Section~\ref{sec:taxonomy}, and the practitioner defines their own types against them.

Crypto-native assets are out of scope. A crypto-native asset is a natively issued crypto asset, Bitcoin and Ether among them, that represents no off-chain asset or liability, whether directly or through a derivative. The definition separates a tokenised dollar and natively issued company stock from the crypto-native class. Many assets commonly described as ``crypto-native'', a fiat-backed stablecoin among them, have off-chain assets and liabilities and so meet the definition of a tokenised asset. Whether the framework evaluates one still turns on the fund correspondence above.

The fund share has the exit the framework charges for: redemption with the issuer at the published valuation over a contractual window. An extension to non-fund assets would rest on identifying how each asset settles. Settlement is either with an issuer or primary dealer under contracted redemption or trading terms, or in a market the practitioner can show to be deep from public data. Justifying liquidity without public data would take a strong evidence standard.

Bridge risk sits outside the scope of this framework. An asset whose token, administration, or valuation oracle depends on a bridge or cross-chain messaging system is not evaluated under the framework.

Feeder structures are eligible. A feeder token holds its economic exposure through a master fund. The look-through of Section~\ref{sec:tax-lookthrough} extends to the master's holdings, read from the master's own documents, so the sensitivities the margin is computed from are the master's. Three master-level risks take named layers (Section~\ref{sec:theory:layers}). Gate pass-through takes the gating layer, and master valuation staleness takes the valuation-staleness layer. The additional administrator in the chain takes the manager/operational layer. The practitioner sets the evidence standard the master-level legs are held to (free parameter, Appendix~\ref{sec:free-parameters}).

Most tokenised assets trade too thinly to offer instant liquidity in large amounts, and that is the ordinary condition of the assets themselves. Their off-chain equivalents settle over days or weeks, and traditional clearing is built around periods of time rather than instant price discovery. \emph{The framework assumes subscription and redemption at the published valuation only}.\footnote{Transfer restrictions can require any party taking delivery of the collateral, a liquidator included, to hold know-your-customer approval from the issuer.} Any automated-market-maker (AMM), on-chain secondary market, slippage, maximal extractable value (MEV), gas, or atomic-withdrawal route is out of scope and deferred to a later draft.

The cost of clearing a position in stress rises with its size, and the framework charges for it. SIMM's risk weights are calibrated on stressed histories. The concentration risk factor steepens the charge once a position outgrows one dealing window of attested capacity. The clearing discount rescales the stressed move to the whole clearing window. What redemption at the published valuation removes is the modelling of venue depth.

The primary object is the liquidity horizon and the process that settles it, which for a tokenised fund is the issuer's own subscription and redemption. The horizon runs from the breach to cash: the worst-case wait to the next submission deadline plus the attested notice, dealing, and settlement legs (Eq.~\eqref{eq:mpor-eff}). Other tokenised assets have other settlement processes and other horizons, and Section~\ref{sec:theory:setting} sets out the three sources those legs are read from, in a fixed order. Redemption at the published valuation is the settlement the framework treats as guaranteed, in stress as well as in normal conditions.

Holding-period risk remains, and it falls on both sides. A borrower waiting to release collateral and the deployment recovering a position in liquidation wait for the same window, and neither can shorten it. The clearing discount $\delta_{\mathrm{clear}}^{(a)}$ reserves against that wait. Section~\ref{sec:theory:stress} sets it precisely. The framework sets the discount per asset, and stays silent on the clearing mechanism and on how the discount is enforced or paid.

\subsection{Outputs and how they are set}\label{sec:overview-recipe}
All four outputs are computed in a single pass from attested terms and published risk weights. The interest cost depends on the liquidation threshold, and the two resolve together in closed form, so there is no circular dependency (Eq.~\eqref{eq:carry-solve}). The haircut $H^{(a)}$ is the fraction of the published valuation reserved against holding a position in liquidation to cash settlement, built as a four-part sum. The liquidation threshold $\theta^{(a)}$ is what remains of the published valuation after the haircut, scaled by the due-diligence tier multiplier $\kappa^{(a)}\le 1$ ($\kappa^{(a)}=1$ at Prime; Section~\ref{sec:policy-tiers}). The borrow limit $\ell_{\max}^{(a)}$ sits the liquidation buffer $\nu^{(a)}$, 4--8 percentage points, below $\theta^{(a)}$. The buffer is set per listing (Section~\ref{sec:policy-notch}). Both caps are keyed to the issuer's attested redemption capacity $R_{\mathrm{cap}}^{(a)}$. The borrow cap is set so that one dealing window clears the pool's entire debt at the point of liquidation, and the supply cap is an elected multiple $M$ of that same capacity:
\begin{gather*}
H^{(a)} = \mathrm{IM}^{(a)} + \delta_{\mathrm{clear}}^{(a)} + \delta_{\mathrm{carry}}^{(a)} + \Lambda^{(a)},\qquad
\theta^{(a)} = \kappa^{(a)}\bigl(1 - H^{(a)}\bigr),\qquad
\ell_{\max}^{(a)} = \theta^{(a)} - \nu^{(a)},\\
B_{\max}^{(a)} \;\le\; R_{\mathrm{cap}}^{(a)},\qquad
S_{\max}^{(a)} \;\le\; M\,R_{\mathrm{cap}}^{(a)}.
\end{gather*}

The borrow cap is further bounded by the committed settlement level, read from the subscription documents case by case as the default level. Where the documents commit to no measurable level the borrow cap is zero, and the asset is listed and not borrowed against (Section~\ref{sec:policy-caps}).

The practitioner sets the model constants outside the calculation. The supply-cap multiple $M$ is an open model quantity (free parameter, Appendix~\ref{sec:free-parameters}; $M=8$ is illustrative). Section~\ref{sec:theory} states the recipe in full and defines each symbol (Eq.~\eqref{eq:recipe}). The four haircut components are:
\begin{enumerate}[label=(\roman*)]
\item $\mathrm{IM}^{(a)}$: the delta component of the ISDA SIMM v2.6 initial margin per unit of published valuation \cite{isda_simm_v2_6}, taken verbatim at SIMM's native ten-business-day horizon. The rescaling to the governing clearing window happens in $\delta_{\mathrm{clear}}^{(a)}$.
\item $\delta_{\mathrm{clear}}^{(a)}$: the \emph{clearing discount}, the difference between the price-risk charge over the clearing window $\mathrm{LH}^{(a)}$ and the margin over ten business days, set by Eq.~\eqref{eq:m-scaling}. It is a charge where the window runs beyond ten business days, and a credit where contractually enforceable terms, attested or documented, enforce a shorter one. Estimated legs and FRTB fallback windows never earn one. Qualified clearing observations can take the total price-risk charge above the rescaled move, which floors it.
\item $\delta_{\mathrm{carry}}^{(a)}$: the interest cost of holding the debt through the clearing window, resolved in closed form at Eq.~\eqref{eq:carry-solve}.
\item $\Lambda^{(a)}$: the sum of additive layers (gating/suspension, valuation staleness/smoothing, default add-on, oracle pipeline integrity, prepayment, manager/operational, wrong-way, leverage add-on, event risk), charging for the risks SIMM structurally cannot see. Each layer is defined in Section~\ref{sec:theory:layers}. The practitioner sets the model constants used to determine the layer values.
\end{enumerate}

Concentration enters inside $\mathrm{IM}^{(a)}$ through the multiplier $\mathrm{CR}=\max\bigl(1,\sqrt{\mathrm{position}/T}\bigr)$, with the threshold $T$ re-keyed to the attested issuer redemption capacity $R_{\mathrm{cap}}^{(a)}$. $R_{\mathrm{cap}}^{(a)}$ is attested per window, window-normalised where the dealing window differs from the basis on which the capacity is attested (Section~\ref{sec:theory:capacity}). It is attested against the contractual dealing terms only, so a holding inside the asset that could fund redemption faster than the contract requires counts as zero capacity unless the issuer attests a committed fast-redemption capacity. A committed fast-redemption component enters $R_{\mathrm{cap}}^{(a)}$ only when the issuer attests it.

Any unavailable input takes a prescribed conservative floor, never zero (Section~\ref{sec:agg-missing}). Controlled records hold the recomputation schedule. The deployed market configuration is the public source for each active LTV, LLTV, borrow cap, and supply cap. Section~\ref{sec:agg-example} works a treasury money-market fund ($H=2.23\%$) end to end.

The practitioner's due-diligence scorecard governs eligibility and tiering. An eligibility failure rejects the asset outright. Otherwise the assigned tier scales the over-collateralisation cushion by the tier multiplier $\kappa^{(a)}\le 1$ (illustratively $\{1.00,\,0.90,\,0.75\}$ for a Prime/Standard/Restricted structure), and it acts on that cushion alone (Section~\ref{sec:policy-tiers}). The multiplier is set at listing from due-diligence quality and re-scored at the governed frequency. An adverse re-score is applied through the new-borrow limit and the caps, never onto the liquidation threshold of a live position.

\subsection{Emergency actions}\label{sec:overview-emergency}
All emergency authority acts on \emph{new} exposure and capacity, never on the liquidation threshold of live positions. Lowering the liquidation threshold on live collateral exactly when the exit route is slowest, on a gated or suspended asset, is procyclical and self-defeating: it would create the forced-selling spiral this design removes. Five actions are available, in increasing severity. The mildest is the new-borrow LTV cut, a lower $\ell_{\max}^{(a)}$ for new borrows only with live positions keeping their limit. Above it sit the cap cut, which reduces $B_{\max}^{(a)}$ and $S_{\max}^{(a)}$, and the origination rate limit, which slows new borrows per dealing window. The two severest are redemption-only mode, in which no new borrows open and collateral exits only through issuer redemption, and the pause, under which the collateral accepts no new borrows and no new supply. The escalation ladder implementing these actions is set out in Section~\ref{sec:operations}.

Guardrails are valuation-native and observable from day one. The operations companion lists these observables, including rating events. Section~\ref{sec:operations} states the rules they implement. Among the rating events, loss of investment grade on rated securitised collateral forces escalation level 3, and \emph{confirmed} loss forces automatic full redemption at escalation level 4. Full redemption files redemption notices at contractual dealing dates with no acceleration, an orderly exit of the asset that still never lowers the liquidation threshold on a live position (Section~\ref{sec:operations}).

The deployed market configuration shows one liquidation threshold per market. An adverse $\theta^{(a)}$ change acts through new-borrow $\ell_{\max}^{(a)}$, caps, and pauses. The live liquidation threshold stays where it was set.
